%Data institusi

%\input{Dozentdaten}

%\renewcommand{\fachbereich}{Bidang studi}

%\renewcommand{\dozent}{Prof. Dr. . }

%Data ujian

\renewcommand{\klausurgebiet}{ }

\renewcommand{\klausurtyp}{ }

\renewcommand{\klausurdatum}{. 20}

\klausurvorspann {\fachbereich} {\klausurdatum} {\dozent} {\klausurgebiet} {\klausurtyp}

%Data untuk tabel poin berikut

\renewcommand{\aeins}{  3  }

\renewcommand{\azwei}{  3   }

\renewcommand{\adrei}{  6  }

\renewcommand{\avier}{  4  }

\renewcommand{\afuenf}{  5  }

\renewcommand{\asechs}{  0  }

\renewcommand{\asieben}{  5  }

\renewcommand{\aacht}{  6  }

\renewcommand{\aneun}{  6  }

\renewcommand{\azehn}{  0  }

\renewcommand{\aelf}{  6  }

\renewcommand{\azwoelf}{  12  }

\renewcommand{\adreizehn}{  0   }

\renewcommand{\avierzehn}{  9  }

\renewcommand{\afuenfzehn}{  65  }

\renewcommand{\asechzehn}{    }

\renewcommand{\asiebzehn}{    }

\renewcommand{\aachtzehn}{    }

\renewcommand{\aneunzehn}{    }

\renewcommand{\azwanzig}{    }

\renewcommand{\aeinundzwanzig}{    }

\renewcommand{\azweiundzwanzig}{    }

\renewcommand{\adreiundzwanzig}{    }

\renewcommand{\avierundzwanzig}{    }

\renewcommand{\afuenfundzwanzig}{    }

\renewcommand{\asechsundzwanzig}{    }

\punktetabellevierzehn

\klausurnote

\newpage

\setcounter{section}{0}

\inputaufgabegibtloesung
{3}
{

Definisikan istilah-istilah berikut
\zusatzklammer {yang dicetak miring} {} {}.
\aufzaehlungsechs{Suatu
\stichwort {kelengkungan utama} {}
pada suatu permukaan diferensiabel

\mavergleichskette
{\vergleichskette
{ Y
}
{ \subseteq }{ \R^3
}
{  }{
}
{    }{
}
{   }{
}
}
{}{}{}
pada titik

\mavergleichskette
{\vergleichskette
{ P
}
{ \in }{ Y
}
{  }{
}
{    }{
}
{   }{
}
}
{}{}{.}

}{Suatu
\stichwort {difeomorfisme} {}
antara manifold diferensiabel
\mathkor {} {L} {dan} {M} {.}

}{Suatu
\stichwort {vektor singgung} {}
di titik

\mavergleichskette
{\vergleichskette
{  P
}
{ \in }{  M
}
{  }{
}
{    }{
}
{   }{
}
}
{}{}{}
dari manifold diferensiabel $M$.

}{Suatu
\stichwort {operator diferensial orde pertama} {}
pada manifold diferensiabel $M$.

}{Suatu
\stichwort {setengah ruang Euklides} {.}

}{Suatu
\definitionsverweis {penghabisan kompak}{}{}
bagi suatu
\definitionsverweis {ruang topologis}{}{}
$X$.
}

}
{} {}

\inputaufgabegibtloesung
{3}
{

Nyatakan teorema-teorema berikut.
\aufzaehlungdrei{Teorema tentang
\definitionsverweis {ruang singgung}{}{}
dari suatu
\definitionsverweis {submanifold tertutup}{}{}

\mavergleichskette
{\vergleichskette
{  M
}
{ \subseteq }{  N
}
{  }{
}
{    }{
}
{   }{
}
}
{}{}{.}}{Teorema tentang bentuk eksplisit dari bentuk diferensial tarik balik $\varphi^* \omega$ oleh pemetaan diferensiabel
\maabbdisp {\varphi} { U } { V
} {}
\zusatzklammer {dengan
\definitionsverweis {himpunan bagian terbuka}{}{}
\mathkor {} {U \subseteq \R^n} {dan} {V \subseteq \R^m} {,}
yang koordinatnya masing-masing dinyatakan dengan
\mathl{x_1  , \ldots , x_n}{} dan
\mathl{y_1  , \ldots , y_m}{}} {} {}
dari suatu
$k$-\definitionsverweis {bentuk diferensial}{}{}
pada $V$ dengan representasi

\mavergleichskettedisp
{\vergleichskette
{  \omega
}
{ =} {\sum_{I,\,  { \# \left(   I  \right) } = k}  f_I  dy_I
}
{ } {
}
{ } {
}
{ } {
}
}
{}{}{,}
dengan
\maabb {f_I} { V } {\R
} {}
adalah
\definitionsverweis {fungsi}{}{.}}{Teorema tentang karakterisasi kurva geodesik terhadap koneksi Levi-Civita pada manifold Riemann $M$.}

}
{} {}

\inputaufgabegibtloesung
{6 (2+4)}
{

Misalkan

\mavergleichskette
{\vergleichskette
{ \alpha,\beta
}
{ \in }{ \R_+
}
{  }{
}
{    }{
}
{   }{
}
}
{}{}{}
dan misalkan

\mavergleichskettedisp
{\vergleichskette
{ Y
}
{ =} { \left\{ (x,y)\in\R^2  \mid   \alpha x^2+\beta y^2 = 1        \right\}
}
{ } {
}
{ } {
}
{ } {
}
}
{}{}{.}
\aufzaehlungzwei {Tentukan
\definitionsverweis {pemetaan Gauss}{}{}
untuk $Y$.
} { Berikan secara eksplisit suatu
\definitionsverweis {pemetaan invers}{}{}
dari pemetaan Gauss untuk $Y$.
}

}
{} {}

\inputaufgabe
{4}
{

Jelaskan pasangan istilah \anfuehrung{ekstrinsik dan intrinsik}{} dalam konteks manifold.

}
{} {}

\inputaufgabegibtloesung
{5}
{

Buktikan rumus luas permukaan suatu benda putar yang dihasilkan oleh
\definitionsverweis {kurva diferensiabel}{}{}
\maabbeledisp {\gamma} {]a,b[} {\R^2
} { t } {(x(t),y(t))
} {,}
dengan

\mavergleichskette
{\vergleichskette
{ y(t)
}
{ > }{ 0
}
{  }{
}
{    }{
}
{   }{
}
}
{}{}{.}

}
{} {}

\inputaufgabe
{0}
{

}
{} {}

\inputaufgabegibtloesung
{5}
{

Misalkan $X$ adalah suatu
\definitionsverweis {torus}{}{.}
Berikan suatu
\definitionsverweis {pemetaan diferensiabel}{}{}
surjektif
\maabbdisp {\varphi} {\R^2} {X
} {}
sedemikian sehingga
\definitionsverweis {pemetaan singgung}{}{}
\maabbdisp {T_P(\varphi)} { T_P\R^2 } { T_{\varphi(P)}X
} {}
juga surjektif di setiap titik

\mavergleichskette
{\vergleichskette
{ P
}
{ \in }{ \R^2
}
{  }{
}
{    }{
}
{   }{
}
}
{}{}{.}

}
{} {}

\inputaufgabegibtloesung
{6 (2+2+2)}
{

Tinjau grafik $M$ dari pemetaan
\maabbeledisp {\varphi} {\R^2} {\R
} {(u,v)} { u^2+uv-v^3
} {,}
sebagai submanifold tertutup berdimensi dua di $\R^3$, yaitu

\mavergleichskettedisp
{\vergleichskette
{ M
}
{ =} { \left\{ (u,v,u^2+uv-v^3)  \mid  (u,v) \in \R^2         \right\}
}
{ } {
}
{ } {
}
{ } {
}
}
{}{}{}
dengan metrik Riemann yang diinduksi dari $\R^3$. Misalkan
\maabbeledisp {\psi} {\R^2} { M
} {(u,v)} { (u,v,u^2+uv-v^3)
} {,}
adalah difeomorfisme yang bersesuaian.

a) Tentukan diferensial total dari $\psi$ serta vektor-vektor citra
\mathl{T_P(\psi) (e_1)}{} dan
\mathl{T_P(\psi) (e_2)}{} di
\mathl{T_{\psi(P)}M}{.}

b) Untuk setiap titik berbentuk

\mavergleichskette
{\vergleichskette
{ P
}
{ = }{ (u,0)
}
{  }{
}
{    }{
}
{   }{
}
}
{}{}{}
tentukan luas jajaran genjang yang direntang oleh
\mathl{T_P(\psi) (e_1)}{} dan
\mathl{T_P(\psi) (e_2)}{} di
\mathl{T_{\psi(P)}M}{}.

c) Untuk setiap titik berbentuk

\mavergleichskette
{\vergleichskette
{ P
}
{ = }{ (0,v)
}
{  }{
}
{    }{
}
{   }{
}
}
{}{}{}
tentukan luas jajaran genjang yang direntang oleh
\mathl{T_P(\psi) (e_1)}{} dan
\mathl{T_P(\psi) (e_2)}{} di
\mathl{T_{\psi(P)}M}{}.

}
{} {}

\inputaufgabegibtloesung
{6}
{

Buktikan aturan hasil kali untuk bentuk diferensial yang diferensiabel pada suatu himpunan terbuka di $\R^n$.

}
{} {}

\inputaufgabe
{0}
{

}
{} {}

\inputaufgabegibtloesung
{6}
{

Misalkan $M$ adalah suatu
\definitionsverweis {manifold berbatas}{}{}
dan misalkan

\mavergleichskette
{\vergleichskette
{ P
}
{ \in }{ \partial M
}
{  }{
}
{    }{
}
{   }{
}
}
{}{}{}
adalah titik batas. Tunjukkan bahwa sifat-sifat berikut ekuivalen bagi suatu vektor singgung

\mavergleichskette
{\vergleichskette
{ v
}
{ \in }{ T_PM
}
{  }{
}
{    }{
}
{   }{
}
}
{}{}{.}
\aufzaehlungzwei {Terdapat suatu lintasan yang terdiferensialkan secara kontinu
\maabb {\gamma} { [- \epsilon, 0]} { M
} {}
dengan

\mavergleichskette
{\vergleichskette
{ \gamma(0)
}
{ = }{ P
}
{  }{
}
{    }{
}
{   }{
}
}
{}{}{}
dan

\mavergleichskette
{\vergleichskette
{ \gamma'(0)
}
{ = }{ v
}
{  }{
}
{    }{
}
{   }{
}
}
{}{}{.}
} { Pada setiap peta bermodel setengah ruang negatif, $v$ dipetakan oleh pemetaan singgung ke setengah ruang kanan.
}

}
{} {}

\inputaufgabegibtloesung
{12}
{

Buktikan teorema Stokes untuk manifold berbatas.

}
{} {}

\inputaufgabe
{0}
{

}
{} {}

\inputaufgabegibtloesung
{9 (1+1+3+2+2)}
{

Tinjau parametrisasi trigonometrik
\maabbeledisp {\psi} {\R } { S^1
} { t } { \left(  \cos  t   , \,   \sin  t                   \right)
} {}
dari lingkaran satuan. Misalkan

\mavergleichskette
{\vergleichskette
{ c
}
{ \in }{ [0, 2 \pi[
}
{  }{
}
{    }{
}
{   }{
}
}
{}{}{}
dipilih tetap. Pada bundel vektor trivial
\maabbdisp {} {S^1 \times \R^2} { S^1
} {}
diberikan suatu
\definitionsverweis {koneksi linear}{}{}
$\nabla$ sedemikian sehingga, sepanjang $\psi$,
\definitionsverweis {koneksi tarik balik}{}{}
$\psi^* \nabla$ diberikan oleh
\definitionsverweis {simbol Christoffel}{}{}
\zusatzklammer {indeks bagi satu-satunya arah diferensiasi dihilangkan} {} {}

\mavergleichskettedisp
{\vergleichskette
{ \begin{pmatrix} \tilde{\Gamma}_1^1 (t)   &  \tilde{\Gamma}_1^2 (t)   \\ \tilde{\Gamma}_2^1(t)  &  \tilde{\Gamma}_2^2(t)  \end{pmatrix}
}
{ =} { \begin{pmatrix}  0  &   - { \frac{  c  }{  2 \pi   } }   \\  { \frac{  c  }{  2 \pi  } }  &  0  \end{pmatrix}
}
{ } {
}
{ } {
}
{ } {
}
}
{}{}{.}
Untuk suatu titik

\mavergleichskette
{\vergleichskette
{ Q
}
{ = }{ \begin{pmatrix}  a  \\b   \end{pmatrix}
}
{ \in }{ \R^2
}
{    }{
}
{   }{
}
}
{}{}{}
tinjau pemetaan
\maabbeledisp {\Psi_{Q}} {\R} { S^1 \times \R^2
} { t } { \left(  \cos  t   , \,   \sin  t  , \,   M(t)  \begin{pmatrix}  a  \\b   \end{pmatrix}                 \right)
} {}
dengan

\mavergleichskettedisp
{\vergleichskette
{ M(t)
}
{ =} { \begin{pmatrix}  \cos  { \frac{  ct  }{  2 \pi   } }   &   -  \sin  { \frac{  ct  }{  2 \pi   } }    \\  \sin  { \frac{  ct  }{  2 \pi  } }   &  \cos  { \frac{  ct  }{  2 \pi   } }  \end{pmatrix}
}
{ } {
}
{ } {
}
{ } {
}
}
{}{}{.}
\aufzaehlungfuenf{Tentukan $\Psi_Q(0)$.
}{Tentukan $\Psi_Q(2 \pi)$.
}{Tunjukkan bahwa $\Psi_Q$ adalah suatu
\definitionsverweis {pengangkatan horizontal}{}{}
sepanjang $\psi$.
}{Tunjukkan bahwa
\maabbeledisp {} {\R^2} { \R^2
} { Q } { \Psi_Q(2 \pi)
} {,}
merupakan
\definitionsverweis {isometri linear}{}{}
\zusatzklammer {titik pangkal $(1,0)$ tidak dicantumkan di sini} {} {.}
}{Tunjukkan bahwa
\maabbeledisp {} {\Z} { \operatorname{SL}_{  2  } \! \left(  \R  \right)
} { k } { \left(  Q \mapsto \Psi_Q(2k \pi )  \right)
} {,}
merupakan
\definitionsverweis {homomorfisme grup}{}{.}
}

}
{} {}

 [[Kategorie:Latexseite]]
